% Generated draft: requires mathematical and source review.
\documentclass[11pt]{article}
\usepackage[margin=25mm]{geometry}
\usepackage{fontspec}
\setmainfont{Latin Modern Roman}
\usepackage{amsmath,amssymb,mathtools,unicode-math}
\setmathfont{Latin Modern Math}
\usepackage{enumitem,needspace,xurl,hyperref}
\setlength{\parindent}{0pt}
\setlength{\parskip}{5pt}
\setlength{\emergencystretch}{3em}
\newcommand{\N}{\mathbb{N}}
\newcommand{\Z}{\mathbb{Z}}
\newcommand{\Q}{\mathbb{Q}}
\newcommand{\R}{\mathbb{R}}
\newcommand{\C}{\mathbb{C}}
\newcommand{\F}{\mathbb{F}}
\newcommand{\A}{\mathbb{A}}
\newcommand{\PP}{\mathbb{P}}
\newcommand{\E}{\mathbb{E}}
\newcommand{\eps}{\varepsilon}
\newcommand{\dd}{\,\mathrm{d}}
\begin{document}
\section*{UnsolvedMath ID: 30004648}
\textbf{Source problem code:} catalogue entry
\subsection{Definitions and mathematical statement}
Let \(p\) be a prime, let \(\nu\geq 1\) be an integer, set \(q=p^\nu\), and let \(K=\F_q(T)\). Let \(G\) be a finite group. Define
\[
p(G)=\bigcap_{\substack{N\trianglelefteq G\\p\nmid |G/N|}}N,
\]
so that \(G/p(G)\) is the maximal quotient of \(G\) whose order is prime to \(p\). For a finite group \(H\), write \(H^{\mathrm{ab}}=H/[H,H]\) and let \(d(H)\) be the minimum cardinality of a generating set of \(H\), with \(d(\{1\})=0\).

The supplied conjecture asks whether, for every such \(p,\nu,G\),
\[
r_{\F_q(T)}(G)=\max\bigl\{1,d\bigl((G/p(G))^{\mathrm{ab}}\bigr)\bigr\}.
\]
The source describes \(r_K(G)\) only as the minimum number of ramified primes. Its precise definition requires qualifications that the supplied record does not specify.

For example, if primes mean all places of \(K\), then for a specified class \(\mathcal C_G\) of finite Galois extensions \(L/K\) with \(\operatorname{Gal}(L/K)\cong G\), the corresponding minimum would be
\[
\min_{L\in\mathcal C_G}
\#\{v\text{ a place of }K:v\text{ ramifies in }L/K\},
\]
provided \(\mathcal C_G\) is nonempty. This is a conditional interpretation, not a determination of the intended invariant. The record does not say whether admissible extensions must have full constant field exactly \(\F_q\), or whether constant field extensions are allowed. It also does not specify whether the place at infinity is counted, whether primes are counted as closed points over \(\F_q\) or as geometric points after extending constants, or what value is assigned when no admissible extension exists. No restriction to tame ramification is supplied. The phrase ``finite group'' includes the trivial group as written, but no separate convention or exclusion for that boundary case is given. These choices must be resolved before the displayed conjecture becomes a fully specified mathematical statement.
\subsection{Short English statement}
For every prime power \(q=p^\nu\) and finite group \(G\), the supplied conjecture predicts that the least number of ramified primes needed to realize \(G\) over \(\F_q(T)\) is determined by the minimum number of generators of the abelianization of its maximal prime-to-\(p\) quotient, with a lower floor of one. The admissible extensions, prime-counting convention, and boundary cases require clarification.
\subsection{Sources}
\begin{itemize}
\item Supplied local UnsolvedMath registry record, permanent ID \(30004648\), title: \emph{Minimal Ramification over Rational Function Fields}; exact source problem code: \texttt{catalogue entry}. Collection: Oberwolfach Reports Open Problems. The record gives proposed year 2021 and names no proposer.
\item Supplied report citation: \emph{Homotopic and Geometric Galois Theory (2022)}, Oberwolfach Reports, pp.~663--744, DOI \texttt{10.4171/owr/2021/12}: \url{https://doi.org/10.4171/owr/2021/12}. The background describes these reports as documenting mathematical workshops and their open-problem discussions. The supplied literature list labels the same reference \emph{Oberwolfach Report 12/2021}; the citation's 2022 designation and the record's 2021 year information are retained as supplied.
\item Supplied literature reference: \emph{Minimal ramification over finite fields}, \url{https://doi.org/10.1093/imrn/rnac370}.
\item Literature assessment dated exactly \texttt{2026-08-22}: the supplied curation labels the problem partially solved and states that the formula is proved for abelian groups and many symmetric and alternating cases, including broad 2023 refinements, but not for arbitrary finite groups and \(q\). This is supplied dated curation derived from the OWR material, not an independently verified literature assessment. The background explicitly requests independent verification of time-sensitive status and cited claims.
\item The registry supplies the statement label \texttt{corrected\_verified} and reports manual semantic, typographic, and LaTeX review asserting self-containedness and preservation of the target. These are supplied editorial claims; this draft does not independently verify them or resolve the formulation issues above.
\end{itemize}
\end{document}
